\section{Aplicaciones e integración funcional}
\label{sec:aplicaciones}

La Entrega 4 incorpora dos clientes sobre un backend distribuido común. La
aplicación web atiende la operación institucional desde navegador y la aplicación
móvil facilita tareas de consulta y registro desde Android. Ambas reutilizan los
mismos contratos HTTP y atraviesan el API Gateway; por ello, la autorización y
las reglas de negocio no se duplican en cada interfaz. Esta organización concreta
los límites descritos en la Sección~\ref{sec:arquitectura-solucion} y conserva la
independencia entre presentación, servicios y persistencia.

\subsection{Backend distribuido}
\label{subsec:backend-funcional}

El backend se compone de cinco aplicaciones Spring que cooperan mediante HTTP.
La separación permite evolucionar capacidades de autenticación, usuarios,
catálogos y reservas de manera acotada, aunque exige gestionar errores parciales
y contratos entre procesos, una consecuencia característica de la comunicación
entre microservicios \cite{aksakalli2021deployment}.

\texttt{auth-service} autentica credenciales, emite y renueva JWT, mantiene
sesiones de renovación y ejecuta recuperación de contraseña. Para construir la
identidad presentada a los clientes consulta el perfil institucional en
\texttt{usuarios-service}. Este último administra perfiles, estudiantes,
docentes y administradores; expone operaciones de consulta, creación,
actualización y cambio de estado, y persiste sus agregados en CockroachDB. Los
cambios de perfil producen eventos consumidos actualmente por un observador de
registro, aplicación del patrón ya analizado en la
Sección~\ref{sec:patrones-gof}.

\texttt{academico-laboratorios-service} gestiona campus, bloques, pisos,
laboratorios, equipos y tipos de equipo, además de facultades, carreras,
materias, periodos lectivos y horarios académicos. Sus API permiten mantener los
catálogos, consultar laboratorios disponibles y obtener información individual.
La operación de detalle completo coordina laboratorio, ubicación y equipos
mediante \texttt{LaboratorioDetalleFacade}; este es el contrato que utiliza el
flujo QR móvil después de identificar un laboratorio.

\texttt{reservas-solicitudes-service} concentra solicitudes, reservas,
disponibilidad y agenda. Permite crear y actualizar solicitudes, consultar
listados, detalle e historial, poner una solicitud en revisión, aprobarla o
rechazarla, cancelarla y proponer o responder una alternativa de horario y
laboratorio. Para las reservas ofrece listado, detalle, calendario y
transiciones de cancelación, inicio, finalización y no asistencia. Strategy
encapsula la comprobación de disponibilidad y State controla las transiciones,
sin trasladar esas reglas a los clientes. La creación y la aprobación emplean
\texttt{Idempotency-Key} y registros persistentes en CockroachDB para impedir
que una repetición equivalente produzca resultados duplicados.

Finalmente, \texttt{api-gateway} es el punto de entrada de Web y Mobile.
Encamina las familias \texttt{/api/v1/auth}, \texttt{/api/v1/perfiles},
\texttt{/api/v1/laboratorios}, \texttt{/api/v1/solicitudes} y
\texttt{/api/v1/reservas} hacia su servicio propietario. Valida firma, emisor,
vigencia y tipo \emph{access} del JWT antes de admitir una ruta protegida. Los
servicios vuelven a aplicar permisos específicos, de modo que una identidad
ausente o inválida genera 401 y una identidad autenticada sin autoridad genera
403. El Gateway también fija cabeceras de seguridad, integra telemetría y aplica
CORS con una lista de orígenes configurable mediante entorno, no con una IP de
despliegue incorporada al código. La concentración de controles en el perímetro
reduce inconsistencias, pero no reemplaza la seguridad interservicio, como
advierten las guías de seguridad para aplicaciones de microservicios
\cite{chandramouli2020nist}.

\subsection{Autenticación y gestión de sesión}
\label{subsec:autenticacion-funcional}

El inicio de sesión admite nombre de usuario o correo institucional y verifica
la contraseña con BCrypt, configurado con factor de coste 12. Tras autenticar,
\texttt{auth-service} recupera el perfil y construye una respuesta con roles,
permisos, token de acceso de corta duración y token de renovación. Las sesiones
de renovación se guardan en CockroachDB mediante el hash SHA-256 del token; no se
persiste su valor en claro. Una renovación valida que la sesión continúe activa,
revoca el token anterior y crea otro dentro de la misma familia. La rotación
reduce la ventana de reutilización de un refresh token, práctica recomendada para
clientes públicos \cite{lodderstedt2025rfc9700}.

La protección de acceso registra intentos y aplica un bloqueo temporal de quince
minutos al alcanzar cinco fallos consecutivos. Un acceso correcto limpia el
contador. La recuperación de contraseña devuelve un mensaje neutral para no
revelar si el identificador existe, limita solicitudes y utiliza tokens de
restablecimiento con expiración e invalidación. Roles y permisos se incluyen en
la identidad autenticada, pero la decisión final se mantiene en las reglas de
Spring Security de cada servicio.

El cierre de sesión web sigue un flujo completo y verificable:

\begin{enumerate}
  \item el cliente lee el refresh token de la sesión y envía
        \texttt{POST /api/v1/auth/logout};
  \item el Gateway dirige la solicitud a \texttt{auth-service};
  \item el servicio valida el token y revoca de forma idempotente la sesión
        persistida;
  \item el cliente elimina access token, refresh token, expiración y usuario de
        \texttt{sessionStorage}, actualiza el estado React y redirige a
        \texttt{/login}.
\end{enumerate}

La limpieza local se realiza incluso si la revocación remota no puede
completarse, para no dejar la interfaz atrapada en una sesión aparente. Esto
prioriza el cierre en el dispositivo, aunque una interrupción de red impide
confirmar la revocación en el servidor y debe considerarse al evaluar el riesgo
residual.

\subsection{Aplicación web}
\label{subsec:aplicacion-web}

La aplicación web es una SPA construida con React~18.3.1 y TypeScript~5.8 en
modo \texttt{strict}. React Router separa las rutas públicas de inicio y
recuperación de contraseña de las rutas autenticadas. Dentro del área protegida,
\texttt{ProtectedRoute} comprueba roles o permisos antes de presentar reservas,
solicitudes, calendario y creación de solicitudes. Esta comprobación mejora la
experiencia al ocultar acciones no disponibles, pero el backend conserva la
autoridad efectiva y responde 403 cuando corresponde.

El contexto de autenticación guarda access token, refresh token, vencimiento y
perfil en \texttt{sessionStorage}. No se afirma el uso de cookies
\texttt{httpOnly}: los JWT están disponibles al código del cliente y se eliminan
al cerrar la sesión o la pestaña. Tampoco se almacenan credenciales en
\texttt{localStorage}. Este último se utiliza exclusivamente para preferencias
no sensibles de idioma y tema mediante Zustand. Cada solicitud protegida añade
el Bearer token; ante un 401, el cliente coordina una sola renovación en curso,
actualiza ambos tokens y reintenta la operación. Si la renovación falla, limpia
la sesión. Los mensajes distinguen sesión expirada (401) de falta de permisos
(403).

En el dominio funcional, la Web presenta laboratorios y métricas de ocupación,
gestión administrativa de perfiles, listas de solicitudes y reservas, detalle e
historial, calendario semanal, creación y cancelación, revisión, aprobación,
rechazo y propuestas alternativas según permisos. Los formularios consultan
catálogos y disponibilidad antes de enviar operaciones. Las pantallas modelan
estados de carga, vacío y error; en listados críticos se ofrece reintento y los
errores HTTP se traducen a mensajes comprensibles.

La internacionalización utiliza i18next con recursos en español e inglés, y el
tema permite modos claro y oscuro. En desarrollo, Vite redirige las rutas
relativas \texttt{/api}; el artefacto de producción se sirve mediante Nginx,
que mantiene la navegación de la SPA y actúa como proxy hacia el Gateway. Por
ello, el navegador no necesita conocer direcciones internas de microservicios.

\subsection{Flujo web de autenticación}
\label{subsec:flujo-auth-web}

La Figura~\ref{fig:flujo-auth-web} resume el recorrido comprobado en el código.
No representa un protocolo adicional: sintetiza las llamadas HTTP y los estados
descritos anteriormente.

\begin{figure}[ht]
  \centering
  \begin{minipage}{0.94\linewidth}
\begin{verbatim}
Usuario -> Login React -> API Gateway -> auth-service
                                   <- Access + Refresh JWT
        -> sessionStorage -> rutas protegidas -> microservicios
        -> 401 -> /auth/refresh -> rotacion -> reintento
        -> logout -> /auth/logout -> revocacion -> limpieza -> /login
\end{verbatim}
  \end{minipage}
  \caption{Flujo funcional de autenticación de la aplicación web.}
  \label{fig:flujo-auth-web}
\end{figure}

\subsection{Aplicación móvil Android}
\label{subsec:aplicacion-mobile}

El cliente móvil se desarrolló en Kotlin con Jetpack Compose, Navigation
Compose y una organización MVVM por funcionalidad. Compila contra Android
API~34, define \texttt{targetSdk 34} y soporta desde \texttt{minSdk 26}.
ViewModels exponen estados de carga, éxito y error mediante \texttt{StateFlow};
repositorios separan las pantallas de Retrofit, Room y el almacenamiento de
sesión. Esta elección sigue prácticas contemporáneas de desarrollo Android con
Kotlin y Jetpack \cite{spath2022proandroid}.

Retrofit y OkHttp consumen el mismo Gateway que la Web. Un interceptor agrega
el access token y un \texttt{Authenticator} serializa la renovación ante 401,
evita bucles de reintentos y reconstruye la solicitud con el token actualizado.
\texttt{EncryptedAuthStorage} conserva access token, refresh token, expiración,
identidad, roles y permisos mediante \texttt{EncryptedSharedPreferences}, con
claves AES-256 respaldadas por \texttt{MasterKey}. La navegación habilita
acciones según permisos. A diferencia de la Web, el logout móvil actual elimina
la sesión local, pero no invoca el endpoint remoto de revocación; esta diferencia
se registra como limitación.

La aplicación permite listar solicitudes y reservas, consultar sus detalles e
historial, crear solicitudes rápidas, verificar disponibilidad, cancelar,
revisar, aprobar o rechazar y gestionar propuestas cuando el perfil posee los
permisos necesarios. También presenta un calendario. Las operaciones de creación
y aprobación envían claves de idempotencia al backend. Los fallos de red y las
respuestas 401, 403 y otros errores HTTP se convierten en estados diferenciados
para la interfaz.

\subsubsection{Persistencia local y funcionamiento sin conexión}

Room define una base local con dos conjuntos de datos. Los incidentes se guardan
en la tabla \texttt{incidentes} y se observan como un flujo reactivo; su creación
es enteramente local. Las reservas obtenidas remotamente se almacenan en
\texttt{reservas\_cache}. Si el listado o el detalle fallan por conectividad, el
repositorio devuelve los datos disponibles en esa caché e informa que la
actualización remota no se completó.

Este comportamiento proporciona lectura degradada, pero no constituye una
sincronización offline completa. No hay cola persistente de escrituras,
WorkManager ni resolución de conflictos; crear o modificar reservas requiere
conectividad. Asimismo, los incidentes no tienen fuente remota ni proceso de
sincronización. La distinción es importante porque una arquitectura
\emph{offline-first} completa debe coordinar explícitamente fuentes locales y de
red, así como la reconciliación posterior \cite{android2026offlinefirst}.

\subsubsection{QR y capacidades del dispositivo}

El flujo QR constituye una capacidad implementada. La pantalla solicita en
tiempo de ejecución el permiso de cámara, utiliza CameraX para previsualización y
análisis, y aplica ML Kit Barcode Scanning sobre el fotograma más reciente. El
valor leído debe representar un UUID de laboratorio; el ViewModel evita procesar
duplicados y, tras validarlo, consulta
\texttt{/api/v1/laboratorios/\{id\}/detalle-completo} a través del Gateway. La
interfaz diferencia QR inválido, fallo de red y error del servicio. Por tanto,
la lectura de cámara y la consulta posterior están conectadas, aunque dependen de
que el laboratorio exista y el backend esté accesible.

\subsubsection{Incidentes y notificaciones}

La aplicación crea y lista incidentes locales y puede producir una notificación
del dispositivo al registrar uno. Para Android~13 o superior solicita
\texttt{POST\_NOTIFICATIONS}; en versiones anteriores utiliza el canal de
notificación configurado. Esta notificación local no debe confundirse con una
comunicación de servidor.

La integración de Firebase Cloud Messaging abarca más que una preparación del
cliente. Android recibe mensajes, crea notificaciones y conserva las
renovaciones de token hasta registrarlas mediante
\texttt{/api/v1/notificaciones/dispositivos}. Reservas persiste dispositivos por
perfil, dispone de un adaptador Firebase condicionado por
\texttt{FIREBASE\_ENABLED} y emite eventos de solicitud, incidente y
planificación. El token no se escribe en logs y los fallos de entrega no revierten
la operación de negocio.

No obstante, \texttt{google-services.json} no está versionado y el despliegue no
aporta en el repositorio credenciales Firebase reales, como corresponde a
secretos. Tampoco existe evidencia conservada de la cadena completa evento
backend $\rightarrow$ Firebase $\rightarrow$ dispositivo con Play Services
$\rightarrow$ notificación recibida. Por ello FCM está
\textbf{implementado parcialmente y pendiente de validación E2E}; no se declara
operativo extremo a extremo. Además, los disparadores encontrados no incluyen
todos los eventos posibles del dominio, por ejemplo la apertura de asistencia.

\subsection{Integración de clientes y servicios}
\label{subsec:integracion-clientes}

La Figura~\ref{fig:integracion-clientes} muestra la integración compartida. Web
y Mobile mantienen presentaciones y almacenamiento propios, pero reutilizan las
mismas API, el esquema Bearer y las decisiones de autorización del backend.

\begin{figure}[ht]
  \centering
  \begin{minipage}{0.88\linewidth}
\begin{verbatim}
WEB ------\
           > API Gateway --> microservicios --> CockroachDB
MOBILE ---/
\end{verbatim}
  \end{minipage}
  \caption{Integración conceptual de clientes, Gateway y servicios.}
  \label{fig:integracion-clientes}
\end{figure}

Los contratos HTTP impiden que los clientes accedan directamente a bases de
datos o dependan de entidades de persistencia. El JWT transporta identidad y
autoridades; el Gateway valida el token y cada servicio aplica permisos sobre
su dominio. Las consultas internas de Auth y Reservas hacia Usuarios o
Académico utilizan una cabecera de API interna configurable por entorno. No se
reproducen sus valores en el informe ni se incorporan secretos al cliente.

\subsection{Seguridad funcional y limitaciones}
\label{subsec:seguridad-limitaciones-apps}

La seguridad funcional combina contraseñas BCrypt, JWT de acceso, refresh tokens
rotados, sesiones persistidas y revocables, RBAC, respuestas diferenciadas 401 y
403, CORS configurable, cabeceras de seguridad y claves para llamadas internas.
Web conserva credenciales exclusivamente en \texttt{sessionStorage}; Mobile las
cifra mediante Android Security Crypto. Ningún cliente incluye tokens en URL y
el servicio FCM no registra el valor recibido. Estos controles reducen riesgos,
pero no justifican afirmar seguridad absoluta; la literatura identifica la
seguridad de microservicios como un problema transversal de identidad,
infraestructura, comunicación y ciclo de desarrollo
\cite{berardi2022microservicesecurity}.

Permanecen cuatro limitaciones funcionales. Primero, incidentes es un módulo
local sin API ni sincronización remota. Segundo, FCM carece de configuración
Firebase verificable, registro de dispositivo y emisor backend. Tercero, salvo
la lectura de reservas almacenadas y las operaciones locales de incidentes, las
funciones remotas dependen de conectividad y disponibilidad del Gateway. Cuarto,
el logout móvil no revoca todavía la sesión de renovación en el backend. La
calidad y cobertura de estas funciones se evalúan en el capítulo de pruebas; no
se infiere su suficiencia a partir de la mera existencia de pantallas o clases.

No se incorporaron capturas: el repositorio no contiene imágenes actuales y
claramente atribuibles a las interfaces E4. Esta omisión evita reutilizar como
evidencia visual las capturas históricas de infraestructura.
